%% template_technical.tex
%% EuMINe DataBridge Hackathon 2026 — Technical Report Template
%%
%% Fill in every TODO section.
%% Compile with: pdflatex template_technical.tex  (run twice)

\documentclass[10pt, a4paper]{article}
\usepackage{eumine_hackathon}

%% ── Title block ────────────────────────────────────────────────────────────
\title{%
  \textbf{Technical Report}\\[4pt]
  {\large EuMINe DataBridge Hackathon \hackathonyear}
}
\author{
  TODO: Team Name \\
  \small TODO: Institution(s) \\
  \small TODO: Contact email
}
\date{TODO: Date submitted}

%% ───────────────────────────────────────────────────────────────────────────
\begin{document}
\maketitle
\thispagestyle{fancy}

% ── Abstract ──────────────────────────────────────────────────────────────
\begin{abstract}
TODO: One paragraph describing your model, the key ideas that make it
better than the MAGPIE + RandomForest baseline, and the final validation
MAE for both target properties.
\end{abstract}

\tableofcontents
\vspace{1em}

% ── 1. Model Architecture ─────────────────────────────────────────────────
\section{Model Architecture}

\subsection{Overview}
TODO: Describe your model at a high level (e.g.\ graph neural network,
equivariant model, transformer, ensemble, etc.) and explain the main
design choices.

\subsection{Input Representation}
TODO: How are crystal structures represented as input?
Examples:
\begin{itemize}
  \item Composition-based descriptors (MAGPIE, Matminer featurisers)
  \item Graph: atoms as nodes, bonds as edges, edge features (distance,
    periodic image offset)
  \item Point cloud / voxel grid
  \item Pre-trained embedding (e.g.\ GNoME, CHGNet, MACE)
\end{itemize}

\subsection{Architecture Details}
TODO: Provide enough detail to reproduce your model.
\begin{itemize}
  \item Number of layers / parameters
  \item Activation functions
  \item Aggregation / readout function
  \item Output head(s) — one head per property or multi-task?
\end{itemize}

If applicable, include a diagram:
\begin{figure}[h!]
  \centering
  % \includegraphics[width=0.7\linewidth]{figures/architecture.png}
  \caption{TODO: Model architecture diagram.}
  \label{fig:architecture}
\end{figure}

% ── 2. Training Procedure ─────────────────────────────────────────────────
\section{Training Procedure}

\subsection{Data Split Used}
TODO: Confirm which split was used (train 700 / val 150 / test 150).
Were any additional data sources incorporated?

\subsection{Loss Function}
TODO: Describe the loss (e.g.\ MAE, MSE, Huber). If multi-task, describe
how per-property losses are combined.

$$\mathcal{L} = \lambda_1 \, \text{MAE}(E_f) + \lambda_2 \, \text{MAE}(E_g)$$

where \(\lambda_1 = \text{TODO}\) and \(\lambda_2 = \text{TODO}\).

\subsection{Optimiser and Hyperparameters}
\begin{table}[h!]
\centering
\caption{Training hyperparameters.}
\label{tab:hparams}
\begin{tabular}{@{}ll@{}}
\toprule
Hyperparameter & Value \\
\midrule
Optimiser           & TODO (e.g.\ Adam) \\
Learning rate       & TODO \\
Scheduler           & TODO \\
Batch size          & TODO \\
Epochs              & TODO \\
Early stopping      & TODO \\
Weight decay        & TODO \\
Random seed         & TODO \\
\bottomrule
\end{tabular}
\end{table}

\subsection{Hardware and Runtime}
TODO: GPU/CPU used, total training time.

% ── 3. Results ────────────────────────────────────────────────────────────
\section{Results}

\begin{table}[h!]
\centering
\caption{Validation-set performance vs.\ baseline.}
\label{tab:results}
\begin{tabular}{@{}lcc@{}}
\toprule
Model & MAE \(E_f\) (eV/atom) & MAE \(E_g\) (eV) \\
\midrule
MAGPIE + RandomForest (baseline) & \(0.XXX\) & \(0.XXX\) \\
\textbf{TODO: Your model name}   & TODO      & TODO      \\
\bottomrule
\end{tabular}
\end{table}

\subsection{Error Analysis}
TODO: Discuss where your model performs well and where it fails.
Include parity plots if space allows:

\begin{figure}[h!]
  \centering
  % \includegraphics[width=0.45\linewidth]{figures/parity_ef.png}
  % \includegraphics[width=0.45\linewidth]{figures/parity_bg.png}
  \caption{TODO: Parity plots for (left) \(E_f\) and (right) \(E_g\)
    on the validation set.}
  \label{fig:parity}
\end{figure}

\subsection{Ablation Studies (optional)}
TODO: Report the effect of removing key components (features, layers, data
sources) to justify your design choices.

% ── 4. Reproducibility ────────────────────────────────────────────────────
\section{Reproducibility}

\subsection{Code Repository}
TODO: Provide a link to your GitHub / GitLab repository.

\subsection{Dependencies}
TODO: Python version and key library versions (e.g.\ PyTorch, pymatgen,
matminer, scikit-learn).

\subsection{Reproducing the Results}
TODO: Step-by-step commands to reproduce training and evaluation:

\begin{lstlisting}[language=bash]
# Example
pip install -r requirements.txt
python train.py --config config/best.yaml --seed 42
python predict.py --model checkpoints/best.pt --out predictions_test.json
\end{lstlisting}

% ── 5. Discussion ─────────────────────────────────────────────────────────
\section{Discussion}

\subsection{Key Findings}
TODO: What did you learn about the MP\,--\,JARVIS dataset split?
Did the cross-source labels cause any issues?

\subsection{Limitations}
TODO: Known failure modes, edge cases, or assumptions that may not hold
outside this dataset.

\subsection{Future Work}
TODO: What would you do with more time / data / compute?

% ── References ────────────────────────────────────────────────────────────
\section*{References}

\begin{enumerate}
  \item Jain et al., \textit{APL Mater.} \textbf{1}, 011002 (2013) —
    Materials Project.
  \item Choudhary et al., \textit{npj Comput. Mater.} \textbf{6}, 173 (2020)
    — JARVIS-DFT.
  \item TODO: Add model-specific references (e.g.\ CGCNN, MEGNet, SchNet,
    DimeNet, GNoME, CHGNet, MACE, \ldots).
\end{enumerate}

\end{document}
